                   11th International Mathematical Competition for University Students
                                         Skopje, 25–26 July 2004

                                          Solutions for problems on Day 1


Problem 1. Let S be an infinite set of real numbers such that |s1 + s2 + · · · + sk | < 1 for every finite subset
{s1 , s2 , . . . , sk } ⊂ S. Show that S is countable. [20 points]
Solution. Let Sn = S ∩ ( n1 , ∞) for any integer n > 0. It follows from the inequality that |Sn | < n. Similarly, if we
define S−n = S ∩ (−∞, − n1 ), then |S−n | < n. Any nonzeroSx ∈ S is an element of some Sn or S−n , because there
exists an n such that x > n1 , or x < − n1 . Then S ⊂ {0} ∪     (Sn ∪ S−n ), S is a countable union of finite sets, and
                                                                     n∈N
hence countable.

Problem 2. Let P (x) = x2 − 1. How many distinct real solutions does the following equation have:

                                       P (P (. . . (P (x)) . . . )) = 0 ?       [20 points]
                                       |    {z       }
                                            2004


Solution. Put Pn (x) = P (P (...(P (x))...)). As P1 (x) ≥ −1, for each x ∈ R, it must be that Pn+1 (x) = P1 (Pn (x)) ≥
                       | {z }
                              n
−1, for each n ∈ N and each x ∈ R. Therefore the equation Pn (x) = a, where a < −1 has no real solutions.
Let us prove that the equation Pn (x) = a, where a > 0, has exactly two distinct real solutions. To this end we
use mathematical induction by n. If n = 1 the assertion follows directly. Assuming that the assertion holds for a
n ∈ N we prove √ that it must also hold
                                    √ for n + 1. Since Pn+1 (x) =√a is equivalent√     to P1 (Pn (x)) = a, we conclude
that Pn (x) = a + 1 or Pn (x) = − a + 1. The equation Pn (x) = a + 1, as       √ a + 1 > 1, has exactly two distinct
real
  √  solutions by  the inductive hypothesis, while  the  equation P n (x) =  −   a + 1 has no real solutions (because
− a + 1 < −1). Hence the equation Pn+1 (x) = a, has exactly two distinct real solutions.
    Let us prove now that the equation Pn (x) = 0 has exactly n + 1 distinct real solutions. Again we √            use
mathematical induction. If n = 1 the solutions are x = ±1, and if n = 2 the solutions are x = 0 and x = ± 2,
so in both cases the number of solutions is equal to n + 1. Suppose that the assertion holds for some n ∈ N .
Note that Pn+2 (x) = P2 (Pn (x)) = Pn2 (x)(Pn2 (x) − 2), so the set of all real solutions of the
                                                                                               √ equation Pn+2 = √ 0 is
exactly the union of the sets of all real solutions of the equations Pn (x) = 0, Pn (x) = 2 and Pn (x) = − 2.
By the inductive
         √         hypothesis√the equation Pn (x) = 0 has exactly n + 1 distinct real solutions, while the equations
Pn (x) = 2 and Pn (x) = − 2 have two and no distinct real solutions, respectively. Hence, the sets above being
pairwise disjoint, the equation Pn+2 (x) = 0 has exactly n + 3 distinct real solutions. Thus we have proved that,
for each n ∈ N , the equation Pn (x) = 0 has exactly n + 1 distinct real solutions, so the answer to the question
posed in this problem is 2005.

                                                   n
                                                         xk , where n ≥ 2, 0 ≤ x1 , x2 , . . . , xn ≤ π2 and
                                                   P
Problem 3. Let Sn be the set of all sums
                                                   k=1

                                                           n
                                                           X
                                                                 sin xk = 1 .
                                                           k=1

   a) Show that Sn is an interval. [10 points]
   b) Let ln be the length of Sn . Find lim ln . [10 points]
                                            n→∞

Solution. (a) Equivalently, we consider the set

                     Y = {y = (y1 , y2 , ..., yn )| 0 ≤ y1 , y2 , ..., yn ≤ 1, y1 + y2 + ... + yn = 1} ⊂ Rn

and the image f (Y ) of Y under

                                      f (y) = arcsin y1 + arcsin y2 + ... + arcsin yn .

Note that f (Y ) = Sn . Since Y is a connected subspace of Rn and f is a continuous function, the image f (Y ) is
also connected, and we know that the only connected subspaces of R are intervals. Thus Sn is an interval.
    (b) We prove that
                                                1                            π
                                                  ≤ x1 + x2 + ... + xn ≤ .
                                              n arcsin
                                                n                            2
Since the graph of sin x is concave down for x ∈ [0, π2 ], the chord joining the points (0, 0) and ( π2 , 1) lies below the
graph. Hence
                                            2x                          π
                                               ≤ sin x for all x ∈ [0, ]
                                            π                            2
and we can deduce the right-hand side of the claim:
                                  2
                                    (x1 + x2 + ... + xn ) ≤ sin x1 + sin x2 + ... + sin xn = 1.
                                  π
The value 1 can be reached choosing x1 = π2 and x2 = · · · = xn = 0.
   The left-hand side follows immediately from Jensen’s inequality, since sin x is concave down for x ∈ [0, π2 ] and
0 ≤ x1 +x2 +...+x
           n
                  n
                    < π2
                            1    sin x1 + sin x2 + ... + sin xn       x1 + x2 + ... + xn
                               =                                ≤ sin                    .
                            n                  n                              n
Equality holds if x1 = · · · = xn = arcsin n1 .
   Now we have computed the minimum and maximum of interval Sn ; we can conclude that Sn = [n arcsin n1 , π2 ].
Thus ln = π2 − n arcsin n1 and
                                                π     arcsin(1/n)  π
                                       lim ln = − lim             = − 1.
                                      n→∞       2 n→∞     1/n      2


Problem 4. Suppose n ≥ 4 and let M be a finite set of n points in R3 , no four of which lie in a plane. Assume
that the points can be coloured black or white so that any sphere which intersects M in at least four points has
the property that exactly half of the points in the intersection of M and the sphere are white. Prove that all of
the points in M lie on one sphere. [20 points]
                                               
                                                 −1, if X is white                                P
Solution. Define f : M → {−1, 1}, f (X) =                           . The given condition becomes X∈S f (X) = 0
                                                  1, if X is black
for any sphere S which passes through at least 4 points of M . For any 3 given points A, B, C in M , denote by
S (A, B, C) the set of all spheres which pass through
                                                 P           Pand at least one other point of M and by |S (A, B, C)|
                                                       A, B, C
the number of these spheres. Also, denote by        the sum X∈M f (X).
    We have                  X      X                                                         X
                    0=                  f (X) = (|S (A, B, C)| − 1) (f (A) + f (B) + f (C)) +                    (1)
                            S∈S(A,B,C) X∈S

since the values of A, B, C appear |S (A, B, C)| times each and the other values appear only once.
    If there are 3 points A, B, C such that |S (A, B, C)| = 1, the proof is finished.        P
    If |S (A, B, C)|P> 1 for any distinct points A, B, C in M , we will prove at first that     = 0.
                                                                                                           n
                                                                                                             
    Assume that        > 0. From (1) itfollows   that f (A) + f (B) + f (C)  <  0 and   summing    by all 3   possible
choices of (A, B, C) we obtain that n3
                                           P                     P
                                             P< 0, which  means     < 0 (contradicts  the  starting assumption).   The
same reasoningP   is applied when assuming      < 0.
    Now, from        = 0 and (1), it follows that f (A) + f (B) + f (C) = 0 for any distinct points A, B, C in M .
Taking another point D ∈ M , the following equalities take place

                                                   f (A) + f (B) + f (C) = 0
                                                   f (A) + f (B) + f (D) = 0
                                                   f (A) + f (C) + f (D) = 0
                                                   f (B) + f (C) + f (D) = 0

which easily leads to f (A) = f (B) = f (C) = f (D) = 0, which contradicts the definition of f .

Problem 5. Let X be a set of 2k−4
                                      
                                   k−2 + 1 real numbers, k ≥ 2. Prove that there exists a monotone sequence
{xi }ki=1 ⊆ X such that
                                             |xi+1 − x1 | ≥ 2|xi − x1 |
for all i = 2, . . . , k − 1.   [20 points]
    Solution. We prove a more general statement:
    Lemma. Let k, l ≥ 2, let X be a set of k+l−4
                                                  
                                            k−2 + 1 real numbers. Then either X contains an increasing sequence
{xi }ki=1 ⊆ X of length k and
                                    |xi+1 − x1 | ≥ 2|xi − x1 | ∀i = 2, . . . , k − 1,
or X contains a decreasing sequence {xi }li=1 ⊆ X of length l and

                                              |xi+1 − x1 | ≥ 2|xi − x1 | ∀i = 2, . . . , l − 1.

   Proof of the lemma. We use induction on k + l. In case k = 2 or l = 2 the lemma is obviously true.
   Now let us make the induction step. Let m be the minimal element of X, M be its maximal element. Let
                                                         m+M                                      m+M
                                Xm = {x ∈ X : x ≤            },         XM = {x ∈ X : x >             }.
                                                          2                                        2
        k+l−4           k+(l−1)−4
                                   + (k−1)+l−4
                                              
Since    k−2        =      k−2        (k−1)−2
                                                 , we can see that either
                                                                                                    
                                           (k − 1) + l − 4                             k + (l − 1) − 4
                             |Xm | ≥                        + 1,      or |XM | ≥                         + 1.
                                             (k − 1) − 2                                    k−2

   In the first case we apply the inductive assumption to Xm and either obtain a decreasing sequence of length l
with the required properties (in this case the inductive step is made), or obtain an increasing sequence {xi }k−1
                                                                                                                i=1 ⊆
Xm of length k − 1. Then we note that the sequence {x1 , x2 , . . . , xk−1 , M } ⊆ X has length k and all the required
properties.
   In the case |XM | ≥ k+(l−1)−4
                                  
                            k−2     + 1 the inductive step is made in a similar way. Thus the lemma is proved.
   The reader may check that the number k+l−4
                                                   
                                               k−2   + 1 cannot be smaller in the lemma.

Problem 6. For every complex number z ∈
                                      / {0, 1} define
                                                  X
                                         f (z) :=     (log z)−4 ,

where the sum is over all branches of the complex logarithm.
   a) Show that there are two polynomials P and Q such that f (z) = P (z)/Q(z) for all z ∈ C \ {0, 1}.             [10
points]
   b) Show that for all z ∈ C \ {0, 1}

                                                             z 2 + 4z + 1
                                                 f (z) = z                .    [10 points]
                                                              6(z − 1)4


    Solution 1. It is clear that
                              P the   left hand side is well defined and independent of the order of summation, because
                                  −4
we have a sum of the type n , and the branches of the logarithms do not matter because all branches are taken.
It is easy to check that the convergence is locally uniform on C \ {0, 1}; therefore, f is a holomorphic function on
the complex plane, except possibly for isolated singularities at 0 and 1. (We omit the detailed estimates here.)
    The function log has its only (simple) zero at z = 1, so f has a quadruple pole at z = 1.
    Now we investigate the behavior near infinity. We have Re(log(z)) = log |z|, hence (with c := log |z|)
                              X                   X                X
                            |    (log z)−4 | ≤        | log z|−4 =    (log |z| + 2πin)−4 + O(1)
                                                  Z ∞
                                              =         (c + 2πix)−4 dx + O(1)
                                                   −∞
                                                      Z ∞
                                                   −4
                                              = c           (1 + 2πix/c)−4 dx + O(1)
                                                         −∞
                                                      Z ∞
                                                   −3
                                              = c           (1 + 2πit)−4 dt + O(1)
                                                               −∞
                                                    ≤ α(log |z|)−3

for a universal constant α. Therefore, the infinite sum tends to 0 as |z| → ∞. In particular, the isolated singularity
at ∞ is not essential, but rather has (at least a single) zero at ∞.
     The remaining singularity is at z = 0. It is readily verified that f (1/z) = f (z) (because log(1/z) = − log(z));
this implies that f has a zero at z = 0.
     We conclude that the infinite sum is holomorphic on C with at most one pole and without an essential singularity
at ∞, so it is a rational function, i.e. we can write f (z) = P (z)/Q(z) for some polynomials P and Q which we
may as well assume coprime. This solves the first part.
     Since f has a quadruple pole at z = 1 and no other poles, we have Q(z) = (z − 1)4 up to a constant factor
which we can as well set equal to 1, and this determines P uniquely. Since f (z) → 0 as z → ∞, the degree of P
is at most 3, and since P (0) = 0, it follows that P (z) = z(az 2 + bz + c) for yet undetermined complex constants
a, b, c.
     There are a number of ways to compute the coefficients a, b, c, which turn out to be a = c = 1/6, b = 2/3.
Since f (z) = f (1/z), it follows easily that a = c. Moreover, the fact lim (z − 1)4 f (z) = 1 implies a + b + c = 1
                                                                                       z→1
(this fact follows from the observation that at z = 1, all summands cancel pairwise, except the principal branch
which contributes a quadruple pole). Finally, we can calculate
                                                                                      
                                X               X                 X          X               1
                  f (−1) = π −4   n−4 = 2π −4       n−4 = 2π −4     n−4 −         n−4  =     .
                                                                                            48
                                nodd                n≥1odd                      n≥1               n≥1even

This implies a − b + c = −1/3. These three equations easily yield a, b, c.
    Moreover, the function f satisfies f (z) + f (−z) = 16f (z 2 ): this follows because the branches of log(z 2 ) =
log((−z)2 ) are the numbers 2 log(z) and 2 log(−z). This observation supplies the two equations b = 4a and a = c,
which can be used instead of some of the
                                       P considerations       above.
                                              1                                      dz
    Another way is to compute g(z) =       (log z)2
                                                    first. In the same way, g(z) = (z−1) 2 . The unknown coefficient d
                                  2
can be computed from lim (z − 1) g(z) = 1; it is d = 1. Then the exponent 2 in the denominator can be increased
                        z→1
by taking derivatives (see Solution 2). Similarly, one can start with exponent 3 directly.
   A more straightforward, though tedious way to find the constants is computing the first four terms of the
Laurent series of f around z = 1. For that branch of the logarithm which vanishes at 1, for all |w| < 12 we have
                                                               w2 w3 w4
                                       log(1 + w) = w −          +   −   + O(|w|5 );
                                                               2   3   4
after some computation, one can obtain
                                  1                    7     1
                                         = w−4 + 2w−2 + w−2 + w−1 + O(1).
                             log(1 + w)4               6     6
The remaining branches of logarithm give a bounded function. So
                                                            7   1
                                  f (1 + w) = w−4 + 2w−2 + w−2 + w−1
                                                            6   6
(the remainder vanishes) and
                                      1 + 2(z − 1) + 76 (z − 1)2 + 16 (z − 1)3   z(z 2 + 4z + 1)
                            f (z) =                                            =                 .
                                                    (z − 1)4                       6(z − 1)4
   Solution 2. ¿From the well-known series for the cotangent function,
                                                    N
                                                    X          1        i    iw
                                              lim                      = cot
                                             N →∞          w + 2πi · k  2     2
                                                    k=−N

and
                              N                                                         i log z
                              X           1         i    i log z  i e2i· 2 + 1   1  1
                      lim                          = cot         = ·i   i log z = +    .
                     N →∞          log z + 2πi · k  2       2     2 e2i· 2 − 1   2 z−1
                            k=−N
Taking derivatives we obtain                                      0
                                       X      1         1      1            z
                                                    = −z ·+           =           ,
                                           (log z)2     2 z−1            (z − 1)2
                                                                  0
                                       X    1       z       z           z(z + 1)
                                                3
                                                  =− ·           2
                                                                      =
                                         (log z)    2    (z − 1)        2(z − 1)3
and                                                                        0
                                                                                    z(z 2 + 4z + 1)
                                                           
                                   X      1       z            z(z + 1)
                                              4
                                                =− ·                            =                   .
                                       (log z)    3            2(z − 1)3              2(z − 1)4
